% GENERATED FILE. Edit canonical lessons, then run npm run generate:books.
% canonical-source-hash: fnv1a64:290c6f3e5dd7d0ca
% canonical-lessons: ES-C42-sintesis-decirlo-una-vez

\chapter{Synthesis --- Saying It Once}
\label{ch:synthesis-say-it-once}
\hlchaptermodality{70}

\begin{chapteropening}
\textbf{By the end of this chapter:} \emph{I can answer a yes/no question about a thing already named without naming it again.}

Complete ES-C42-sintesis-decirlo-una-vez: answer Tienes el libro and Haces la comida with lo tengo and la hago, and say what was wrong with the grammatically correct version.
\end{chapteropening}

\section[Practice]{Synthesis — saying it once}

\label{lesson:ES-C42-sintesis-decirlo-una-vez}



\begin{quote}
Answer out loud: \emph{¿Tienes el libro?} Say yes, and do not use the word \emph{libro}.
\end{quote}

\subsection*{The exchange}
Here are the same two exchanges, twice. Read both versions aloud.

\textbf{Version one:}

\begin{quote}
— \emph{¿Tienes el libro?} — \emph{Sí, tengo el libro.} — \emph{¿Haces la comida?} — \emph{Sí, hago la comida.}
\end{quote}

\textbf{Version two:}

\begin{quote}
— \emph{¿Tienes el libro?} — \emph{Sí, lo tengo.} — \emph{¿Haces la comida?} — \emph{Sí, la hago.}
\end{quote}

Every sentence in version one is correct. Not one of them is a mistake. And no speaker of Spanish would say them.

\begin{grammarlens}[title={Grammar Lens: what the pronoun is actually for}]
These four words do not make you more \textbf{correct}. They make you sound like somebody having a conversation instead of somebody completing an exercise.

A conversation carries what has already been said. English does the same thing and you have never noticed: you would not answer \emph{do you have the book?} with \emph{yes, I have the book}.

So this is the first place where knowing the form is not the skill. \textbf{The skill is noticing that the thing has already been named} — and no table can drill that.
\end{grammarlens}

\subsection*{Your turn}
Say these aloud:

\begin{itemize}
\raggedright
  \item both versions of the book exchange, and hear the difference
  \item \textquotedblleft{}Sí, lo tengo\textquotedblright{} and \textquotedblleft{}Sí, la hago\textquotedblright{} until no decision is needed
\end{itemize}

\begin{tcolorbox}[breakable,colback=teal!4,colframe=teal!35!black,title={Before you move on}]
\emph{¿Tienes el libro?} (\emph{\textbf{Sí, lo tengo}}.) \emph{¿Haces la comida?} (\emph{\textbf{Sí, la hago}}.) And what were these four words for? (\textbf{Not} repeating what has already been said.) Next: the past tense — and you will find these four words sit in exactly the same place there.
\end{tcolorbox}
